




Recall that a Poincaré complex has a visible signature, originally due to Weiss-Williams, and described in the language of Poincaré $\infty$-categories in \S\refone{subsection:parametrised-spectra}. The purpose of this appendix is to give a second perspective on this visible signature. 
knowledge of either or the two mentioned sources, except for the very last part where we compare it to the construction from loc.\ cit.

Let us give an outline of the approach. Recall that the $\infty$-category of finite parametrised spectra over a space $X$, in which the underlying object of the visible signature lies, models the colimit of the constant $X$-indexed diagram with value $\Spaf$. This is an instance of a general phenomenon: Given a diagram 
\[
F\colon X \to \Catx
\]
of stable $\infty$-categories, the collection of Ind-right adjoints $i_x^R\colon \colim_X F \to \Ind(F(x))$ of the colimit structure maps $i_x\colon F(x) \to \colim_X F$ induce a fully-faithful embedding $e_{\Ind}\colon \colim_X F \to \lim_X \Ind\circ F$. Now when $X$ is finite the essential image of $e_{\Ind}$ contains $\lim_X F$. 
The resulting functor $\lim_X F \to \colim_X F$ is an instance of a rather general limit to colimit norm map, which we call the cocartesian conorm (although we refrain here from describing it in the most general setting). In the case of parametrised spectra, this corresponds to the inclusion of pointwise finite parametrised spectra inside of finite parametrised spectra. Dually, instead of taking the Ind-right adjoints $i_x^R$ one can take the Pro-left adjoints $i_x^L\colon \colim_X F \to \Pro(F(x))$, yielding an embedding $e_{\Pro}\colon \colim_X F \to \lim_X \Pro\circ F$, which again contains the image of $\lim_X F$ when $X$ is finite. This gives a dual variant $\lim_X F \to \colim_X F$ of the conorm map, which we call the cartesian conorm. Returning to the case where $F$ is constant with value $\Spaf$, it turns out that when $X$ is a Poincaré complex the cartesian and cocartesian conorm maps differ by tensoring with the dualising complex $\nu\colon X \to \Pic(\Spaf)$, and in particular share the same image. This common image is then invariant 
under the duality associated to the Poincaré structure of $\colim_{x \in X} (\Spaf,\QF^{\uni}\otimes \nu_x)$ (which is where the visible signature lives as a Poincaré object). In addition, restricted along the cartesian norm map, the resulting Poincaré structure on $\lim_X \Spaf$ identifies it with the limit $\lim_X (\Spaf,\QF^{\uni})$ of the constant diagram with value $(\Spaf,\QF^{\uni})$. In particular, this full Poincaré subcategory contains a canonical Poincaré object, corresponding to the constant diagram with value to the universal Poincaré object $(\SS,q_{\uni})$. This is our alternative construction of the visible signature.  




\subsection{The (co)cartesian conorm}

Let $X$ be a finite space and $F\colon X\to \Catx$ a functor. The cocartesian conorm is an exact functor
\[
\ccconorm\colon \lim_X F \to \colim_X F
\]
defined as follows: The natural transformations of $X$-indexed diagrams %
\[
\lim_X F \xrightarrow{p_x} F(x) \xrightarrow{i_x} \colim_X F
\]
compose to a natural transformation between constant diagrams, which we may interpret as a functor 
\[
X \to \Hom_{\Catx}(\lim_X F, \colim_X F) \subset \Funx(\lim_X F, \colim_X F).
\]
The cocartesian conorm map is then, by definition, obtained by taking the colimit of this $X$-indexed family.


We emphasise that the colimit $\colim_X F$ is to be taken in $\Catx$, which (unlike the limit) does not agree with the colimit in $\Cat$. An explicit description of the colimit may be obtain as follows. For each $x \in X$ the associated exact functor $i_x\colon F(x) \to \colim_X F$ admits an Ind-right adjoint $i_x^R\colon \colim_X F \to \Ind(F(x))$, which together determine an exact functor
\[
e_{\Ind}\colon \colim_X F \to \lim_{X\op} \Ind\circ F = \lim_X\Ind \circ F ,
\]
where we have used the canonical equivalence $X \simeq X\op$.
The functor $e_{\Ind}$ is fully-faithful and exhibits its target as the Ind-completion of its source. To see this, note that $\Ind$ is colimit preserving when viewed as a functor from $\Catx$ to $\Pr^L_{\mathrm{st}}$ (as follows from the universal mapping property $\Fun^L(\Ind(\C),\D) = \Funx(\C,\D)$, see~\cite{HTT}*{Proposition 5.3.6.2, Example 5.3.6.8}). It follows that the collection of functors $\Ind(i_x)$ exhibit $\Ind(\colim_X F)$ as the colimit of $\Ind\circ F$ in $\Pr^L_{\mathrm{st}}$, and hence the collection of their right adjoints $i_x^R$ exhibit $\Ind(\colim_X F)$ as the limit of $\Ind \circ F$, but viewed as a diagram $X=X\op \to \Pr^R_{\mathrm{st}}$. This limit is then calculated in large categories %
(see~\cite{HTT}*{Theorem 5.5.3.18}), or in large stable ones for that matter,
and so the claim follows. This leads to a concrete description of $\colim_X F$ as the stable subcategory of 
$\lim_{X} \Ind\circ F$
spanned by the images of the composites 
\[
F(x) \to \Ind F(x) \xrightarrow{x_!} \lim_{X} \Ind\circ F
\]
where $x_!$ is left adjoint to the projection $p_x\colon \lim_{X} \Ind\circ F \to \Ind F(x)$ from the limit. %
Under this description we can describe the cocartesian conorm as follows:

\begin{lemma}
\label{lem:cocartesian_norm}%
For any finite space $X$ and $F\colon X\to \Catx$ the cocartesian conorm and the functor $e_{\Ind}$ above fit into a commutative triangle 
 \begin{equation}
\begin{tikzcd}
 & \lim_X F \ar[ld, "\ccconorm"'] \ar[d, "\mathrm{incl}"] \\
 \colim_X F \ar[r, "e_{\Ind}"'] & \lim_X \Ind\circ F.
\end{tikzcd}
\end{equation}
with vertical map the canonical inclusion, and all (co)limits computed in (possibly large) stable $\infty$-categories. In particular, the cocartesian conorm is fully faithful.
\end{lemma}
\begin{proof}
By abuse of notation let us write $p_x$ for both projections $\lim_X F \to F(x)$ and $\lim_X \Ind\circ F \to \Ind(F(x))$ and let $x_!$ be the left adjoint of the latter as above. %
We then calculate
\[
e_{\Ind}(\ccconorm(\vphi)) = e_{\Ind}(\colim_{x \in X}i_xp_x(\vphi)) = \colim_{x \in X}e_{\Ind}i_xp_x(\vphi) = \colim_{x \in X} x_!p_x(\vphi) \xrightarrow{\simeq} \vphi ,
\]
where the last equivalence, induced by the collection of counit maps $x_!p_x(\vphi) \to \vphi$, holds more generally for any $\vphi \in \lim_X \Ind\circ F$, as can be verified by mapping into an arbitrary test object $\psi \in \lim_X \Ind\circ F$.
\end{proof}



The \defi{cartesian conorm} 
\[
\cconorm\colon \lim_X F\to \colim_X F
\]
is the opposite of the cocartesian conorm map applied to $\mathrm{op}\circ F$, after identifying 
\[
\lim_X \mathrm{op}\circ F \simeq (\lim_X F)\op \quad \mathrm{and} \quad \colim_X \mathrm{op}\circ F \simeq (\colim_X F)\op.
\]
Equivalently, it is given by replacing in the formula the colimit by a limit, that is, $\cconorm(\vphi) = \lim_x i_xp_x(\vphi)$. %
A dual version of Lemma~\reffour{lem:cocartesian_norm} can then be obtained by replacing $e_{\Ind}$ with the dual construction
\[
e_{\Pro}\colon \colim_X F \to \lim_{X} \Pro\circ F
\]
induced by the $X$-indexed family of Pro-left adjoints $i_x^L\colon \colim_X F \to \Pro(F(x))$.
The functor $e_{\Pro}$ is then fully faithful and exhibits the target as the $\Pro$-completion of the source. As in Lemma \reffour{lem:cocartesian_norm} above, the cartesian conorm and the embedding $e_{\Pro}$ then fit in a commutative triangle
\begin{equation}
\label{eq:cartesian_conorm_and_e_Pro}%
\begin{tikzcd}
 & \lim_X F \ar[ld, "\cconorm"'] \ar[d, "\mathrm{incl}"] \\
 \colim_X F \ar[r, "e_{\Pro}"'] & \lim_X \Pro\circ F.
\end{tikzcd}
\end{equation}
with vertical map the canonical inclusion. In particular, the cartesian conorm is fully faithful.





The following result explains how to interpret the cartesian conorm in the (arguably more common) embedding $e_{\Ind}$. The formula for this uses the ``fibrewise'' tensoring of $\lim_X G$, for $G\colon X\to \PrL_{\mathrm{st}}$, over $\Spa^X$, the category of parametrised spectra over $X$, 
\[
\otimes \colon (\lim_X G) \otimes \Spa^X \to \lim_X G
\]
defined through the $X$-family of exact functors
\[
(\lim_X G)\otimes \Spa^X \xrightarrow{x^*\otimes x^*} G(x)\otimes \Spa \to G(x)
\]
where the last map is the usual tensoring of the presentable stable $\infty$-category $G(x)$ over spectra.

For what follows, recall that the dualising complex $\nu_X$ of a space $X$ is the universal parametrised spectrum equipped with a natural transformation 
\[
T\colon \colim_{x\in X} \xi\otimes \nu_X \to \lim_{x\in X} \xi
\]
for $\xi\in \Spa^X$. This transformation is %
an equivalence whenever $X$ is compact (in particular, finite). 

\begin{lemma}
\label{lem:e_Ind_of_cartesian_conorm}%
For any finite space $X$ and $F\colon X\to \Catx$ the %
composite functor 
\[
\lim_X F \xrightarrow{\cconorm} \colim_X F \xrightarrow{e_{\Ind}} \lim_X \Ind\circ F
\]
is canonically equivalent to tensoring with the dualising complex of $X$.
\end{lemma}


\begin{remark}
Applying this lemma to the constant functor $F=\const_{\Spa^f}\colon X\to \Cat$, and the object $\const_\SS$ in $\lim_X F = (\Spa^f)^X$,  we find that the dualising complex $\nu_X$ is explicitly given by the formula 
\[
\nu_X \simeq \lim_{x\in X} x_! \SS
\]
which corresponds to the definition of the dualising complex originally given by Klein in \cite{klein-dualizing}.
\end{remark}

For the proof of Lemma~\reffour{lem:e_Ind_of_cartesian_conorm}, we first observe that the transformation $T\colon \colim_{x\in X} \xi\otimes \nu_X \to \lim_{x\in X} \xi
$ %
extends to a similar transformation valid for arbitrary functors $\xi\colon X\to \C$ taking values in compactly generated stable $\infty$-categories (which is again an equivalence whenever $X$ is compact), which can be seen by identifying $\C\simeq \Ind(\C^\omega) \simeq \Funx((\C^\omega)\op, \Spa)$, and using that limits and colimits in $\Funx((\C^\omega)\op, \Spa)$ are computed pointwise.

\begin{proof}[Proof of Lemma \reffour{lem:e_Ind_of_cartesian_conorm}]
The composite functor $e_{\Ind} \circ \cconorm$ is given by 
\[
e_{\Ind}\big(\lim_{x \in X}i_xp_x(\vphi)\big) = \lim_{x\in X} x_!(p_x(\varphi))= \colim_{x\in X} (x_!(p_x(\varphi)) \otimes p_x(\nu_X)) = \colim_{x\in X} x_! p_x(\varphi\otimes \nu_X) = \varphi\otimes \nu_X
\]
where we used the equivalence $T$ in the compactly generated stable $\infty$-category $\lim_X \Ind\circ F$.
\end{proof}



\medskip



\subsection{The visible signature via the conorm map}

We now consider the analogue of the conorm map construction for diagrams of Poincaré $\infty$-categories. 
Suppose given a diagram $F\colon X \to \Catp$ where $X$ is a finite space. As above we may consider the resulting $X$-indexed family of Poincaré functors
\[
i_xp_x\colon \lim_X F \to \colim_X F .
\]
We then encounter the following immediate obstacle: neither of the exact functors $\cconorm = \lim_x i_xp_x$ and $\ccconorm = \colim_xi_xp_x$ is duality preserving; in fact, since each of the individual $i_xp_x$ is duality preserving we have that $\cconorm$ and $\ccconorm$ are the duals of one another in the Poincaré $\infty$-category $\Funx(\lim_X F,\colim_X F)$.
To solve this point, we first restrict to the case where the finite space is a \emph{Poincaré complex}, that is, its dualising complex $\nu_X$ takes values in invertible spectra. This in itself does not remove the above obstacle. However, it allows us to consider a twisted version of the conorm map, in the form a Poincaré functor
\[
\lim_X F \to \colim_X(F\otimes \nu_X) = \colim_{x \in X} F(x) \otimes \nu_x
\]
where for a hermitian $\infty$-category $(\C,\QF)$ and a spectrum $E$ we write $(\C,\QF)\otimes E := (\C,\QF\otimes E)$, which is Poincaré if $(\C,\QF)$ is so and $E$ is invertible. The exact functors $i_xp_x$ are then not duality preserving, but rather become $\nu_x\otimes i_xp_x$ upon conjugating with the dualities on both sides. Conjugating $\lim_x i_xp_x$ by the duality then gives $\colim_x \nu_x\otimes i_xp_x \simeq \lim_x i_xp_x$, so that the cartesian conorm %
has a chance of being promoted to a Poincaré functor. This is indeed the case:


\begin{proposition}
\label{proposition:hermitian_conorm}%
Let $X$ be a finite space with dualising complex $\nu_X$, and $F\colon X\to \Cath$ be a functor. 
The \emph{cartesian} conorm naturally upgrades to a hermitian functor 
 
\[
\cconorm\colon \lim_X F \to \colim_X (F\otimes \nu_X)
\]
which is in fact a hermitian embedding (i.e., identifies the source with a hermitian subcategory of the target).   
Furthermore, if $X$ is a Poincar\'e complex and $F$ takes values in $\Catp$, then this hermitian functor is a Poincar\'e functor between Poincar\'e categories.
\end{proposition}

Before giving a proof of this lemma, we employ it for the promised description of the visible signature. For a Poincar\'e complex $X$ with dualising complex $\nu_X$, we consider the constant functor $X\to \Catp$ with value $(\C, \QF)$. Its limit and colimit are the cotensor $(\C, \QF)^X$ and $(\C, \QF)_X$ respectively. We also write 
\[
(\C, \QF)_{(X, \nu_X)}:= \colim_X ((\C, \QF)\otimes \nu_X).
\]
According to Proposition~\reffour{proposition:hermitian_conorm}, the cartesian conorm thus refines to a Poincar\'e functor 
\[
\cconorm\colon (\C, \QF)^X \to (\C, \QF)_{(X, \nu_X)}.
\]
Applying $\Poinc(-)$ gives us a map of spaces
\[
\cconorm \colon \Hom_{\Sps}(X, \Poinc(\C, \QF)) \simeq \Poinc((\C, \QF)^X) \to \Poinc((\C, \QF)_{(X, \nu_X)}).
\]

\begin{definition}
\label{definition:visible-signature}%
The \defi{visible signature} of a Poincar\'e complex $X$ is the element 
\[
\vsign(X) \in \Poinc((\Spa^f, \QF^{\uni})_{(X, \nu_X)})
\]
obtained by applying this construction to the constant map $X\to \Poinc(\Spa^f, \QF^{\uni})$ with value the universal Poincar\'e object $\SS^u\in \Poinc(\Spa^f, \QF^{\uni})$. 
\end{definition}


We come to the proof of Proposition~\reffour{proposition:hermitian_conorm}. We first recall the $\Pro$-completion of a hermitian category $(\C, \QF)$, denoted $\Pro(\C, \QF)$: This is the stable $\infty$-category $\Pro(\C)$ equipped with the canonical colimit-preserving extension $\hat\QF\colon \Pro(\C)\op\simeq \Ind(\C\op)\to \Spa$, i.e., the hermitian structure left Kan extended from the inclusion $i\colon \C\to \Pro(\C)$. This canonically upgrades $i$ to a hermitian functor (in fact, hermitian embedding), and evidently the whole construction can be made functorial in $(\C, \QF)\in\Cath$. 

\begin{lemma}
\label{lemma:e_Pro_hermitian}%
Let $X$ be any space with dualising complex $\nu_X$, and $F\colon X\to \Cath$. The embedding $e_{\Pro}$ naturally refines to a hermitian functor
\[
e_{\Pro}\colon \colim_X (F\otimes \nu_X) \to \lim_X \Pro\circ F.
\]
If $X$ is compact, this hermitian functor exhibits the target as the $\Pro$-completion of the source hermitian category, and so is, in particular, a hermitian embedding.
\end{lemma}

\begin{proof}
 The hermitian structure on the colimit $\colim_X (F\otimes \nu_X)$ is formed, according to \S\refone{subsection:limits}, by left Kan extending each constituent $\QF_x(-)\otimes \nu_X(x)$ along the opposite of  the inclusion $i_x\colon F(x)\to \colim_X F$, and taking the colimit over these. Thus, its further left Kan extension along the $\Pro$-completion $e_{\Pro}$ (of stable $\infty$-categories) is given by the formula 
\[
(\lim_X \Pro\circ F)\op\to \Spa, \quad \varphi\mapsto \colim_{x\in X} \hat \QF_x(x^* \varphi)\otimes \nu_X(x)
\]
(since the left Kan extension along (the opposite of) $e_{\Pro}\circ i_x = x_*$ is given by pullback along its left adjoint $x^*$). 
On the other hand, again by \S\refone{subsection:limits}, the hermitian structure on $\lim_X \Pro\circ F$ is given by 
\[
(\lim_X \Pro\circ F)\op\to \Spa, \quad \varphi\mapsto \lim_{x\in X} \hat  \QF_x(x^* \varphi).
\]
The claimed refinement is given by the the transformation 
\[
\colim_{x\in X} \hat  \QF(x^*\varphi) \otimes \nu_X(x) \to \lim_{x\in X} \hat  \QF(x^*\varphi)
\]
issued from the defining property of the dualising complex. If $X$ is compact, this transformation is an equivalence. 
\end{proof}



\begin{proof}[Proof of Proposition~\reffour{proposition:hermitian_conorm}]
We consider the diagram in $\Cath$ 
 \begin{equation}
\label{eq:cocartesian_conorm_and_e_Ind}%
\begin{tikzcd}
 & \lim_X F \ar[ld, "\cconorm"', dotted] \ar[d, "\mathrm{incl}"] \\
 \colim_X (F\otimes \nu_X) \ar[r, "e_{\Pro}"'] & \lim_X \Pro\circ F.
\end{tikzcd}
\end{equation}
where the vertical map is the limit of the canonical hermitian embeddings (which is again a hermitian embedding) and the lower line is issued from Lemma~\reffour{lemma:e_Pro_hermitian}. We obtain a unique dotted lift from~\eqreffour{eq:cartesian_conorm_and_e_Pro},
which is again a hermitian embedding. 

It will suffice to show that for any $\varphi\in \lim_X F$, the spectral presheaf
\[
\Bil_{\colim_X (F\otimes \nu_X)}(-, \cconorm(\varphi))
\]
is represented by some object of the form $\cconorm(\psi)$ (where $\psi$ will then necessarily be the dual of $\varphi$ in $\lim_X F$). By the hermitian embedding $e_{\Pro}$ from Lemma~\reffour{lemma:e_Pro_hermitian}, this spectral presheaf is restricted from the spectral presheaf 
\[
\Bil_{\lim_X \Pro\circ F}(-, e_{\Pro}\cconorm \varphi) \simeq \Bil_{\lim_X \Pro\circ F}(-, \varphi) \simeq \lim_x \Bil_{\Pro F(x)}(p_x(-),p_x(\vphi)),
\]
and the latter is represented in $\lim_X \Pro\circ F$ by the pointwise dual $\Dual_{\lim F} \vphi$.
\end{proof}



\medskip



\begin{remark}
We mention the following two variants of this construction:
\begin{enumerate}
\item A version of the (co-)cartesian conorm still exists for spaces $X$ that are merely compact %
if the target category $\colim_X (F\otimes \nu_X)$ is replaced by its idempotent completion (that is, we form the colimit in $\Catxi$). Indeed, any idempotent complete stable $\infty$-category admits limit indexed over compact spaces, and the rest of the proof works verbatim. Consequently, a compact Poincar\'e complex admits a visible signature 
\[
\vsign(X)\in \Poinc(((\Spa^f, \QF^{\uni})_{(X, \nu_X)})^\natural).
\]
Under $e_{\Ind}$, this corresponds to the category of compact parametrised spectra (as opposed to finite ones). 
 \item For a finite space $X$, our argument really shows that the cartesian conorm is a Poincar\'e functor provided each $F(x)$ and each $F(x)\otimes \nu_X(x)$ are Poincar\'e categories.
For instance, if $F$ is constant with value $(\Mod^\omega(R), \QF^\tate)$
for an $\Einf$-ring spectrum $R$, this amounts to saying that each $\nu_X(x)\otimes R$ is invertible in $\Mod{R}$. We deduce that for such an ``$R$-Poincar\'e complex'', we can construct an ``$R$-visible signature''
\[
\vsign(X; R)\in \Poinc((\Mod^\omega(R), \QF^\tate)_{(X, \nu_X)})
\]
by applying the above construction to the constant map $X\to \Poinc(\Mod^\omega(R), \QF^\tate)$ with value the tautological Poincar\'e object $R$. 
\end{enumerate}
\end{remark}


\medskip

We conclude with a comparison of the present approach with the one pursued in \S\refone{subsection:parametrised-spectra}. In loc.\ cit.\ we denote by $\Spaf_X \subseteq \Sp^X$ the full subcategory spanned by the finite parametrised spectra. %
From the present point of view, it is exactly the essential image under $e_{\Ind}$ %
of the colimit of the constant $X$-diagram on $\Spaf$. 
Note the slight discrepancy in conventions: for a space $X$ we generally denote by $\C_X$ the abstract colimit of the constant diagram on $\C$, while in the case of $\Spaf_X$ we implicitly implement the embedding $e_{\Ind}$ and view it as a full subcategory of $\Sp^X$.  
For given spherical fibration $\xi\in \Pic(\SS)^X$, we then defined the \emph{$\xi$-twisted visible hermitian structure} $\QF^\vis_\xi$ on $\Spaf_X$, with linear and bilinear parts given as
\[
 \Lin^\vis_{\xi} := \Hom_{\Spa^X} (-, \xi) \quad \mathrm{and} \quad \Bil^\vis_{\xi} := \Hom_{\Spa^{X\times X}}(-\boxtimes -, \Delta_!(\xi)) 
\]
and obvious connecting transformation $\Lin^\vis_{\xi}\to (\Bil^\vis_{\xi})^\tC$. We write $
\Dual_{\xi}$ for the duality of $\QF^{\vis}_{\xi}$.
This construction is covariantly functorial in $(X,\xi)$ via left Kan extension of parametrised spectra, %
and we show in Proposition~\refone{proposition:visible-is-colimit-full} that the collection of Poincaré functors $x_!\colon (\Spaf,\QF^{\uni})\otimes \xi_x = (\Spaf_{\{x\}},\QF^{\uni}_{\xi_x}) \to (\Spaf,\QF^{\vis}_\xi)$ induces an equivalence
\[
\colim_{x \in X} (\Spaf,\QF^{\uni})\otimes \xi_x \xrightarrow{\simeq} (\Spaf,\QF^{\vis}_\xi)
\]
in $\Catp$. 


Now suppose that $X$ is a Poincaré complex with dualising complex $\nu_X$. Note that the Ind-right adjoint of $x_!\colon \Spaf \to \Spaf_X$ is given by evaluation at $x$, and since $x_!$ is duality preserving, for instance, for the dualities associated to the trivial spherical fibrations $\SS_X$ and $\SS$ on $X$ and $*$, respectively, the Pro-left adjoint of $x_!$ is given by $\vphi \mapsto \hat{\Dual}_{\SS}x^*(\Dual_{\SS_X}\vphi)$, where $\hat{\Dual}_\SS\colon \Ind(\Spaf)\op = \Pro((\Spaf)\op) \to \Pro(\Spaf)$ is the equivalence induced by $\Dual_{\SS}\colon (\Spaf)\op \to \Spaf$. Proposition~\reffour{proposition:hermitian_conorm} then implies that the collection of Pro-left adjoints $\hat{\Dual}_{\SS}x^*(\Dual_{\SS_X}-)\colon \Spaf \to \Pro(\Spaf)$ induces a hermitian embedding
\[
(\Spaf,\QF^{\vis}_{\xi}) \to \lim_X \Pro(\Spaf,\QF^{\vis}_{\xi\otimes\nu_x^{-1}}) ,
\]
given on the level of underlying objects by
\begin{equation}
\label{equation:ind-to-pro}%
\vphi \mapsto \hat{\Dual}^{\fib}\Dual_{\SS_X}(\vphi) 
\end{equation}
where $\hat{\Dual}^{\fib} = \lim_X \hat{\Dual}_{\SS}\colon (\Ind(\Spaf)^X)\op \to \Pro(\Spaf)^X$ is given by applying $\hat{\Dual}_\SS$ fibrewise.
In particular, we obtain an equivalence
\[
\QF^{\vis}_{\xi}(\Dual_{\SS_X}\vphi) \simeq \lim_{x \in X} \QF^{\uni}(\Dual_{\SS_X}\vphi(x))\otimes \xi_x \otimes \nu_x^{-1} \simeq \colim_{x \in X}\QF^{\uni}(\Dual_{\SS}\vphi(x))\otimes \xi_x,
\]
a fact that was also observed in Remark~\refone{remark:useful-formula} (which, as written, also applies when $X$ is not finite, but reduces to the this identification when $X$ is finite). The construction of the visible signature in \S\refone{subsection:parametrised-spectra} then proceeds by noting that $\nu_X = \Dual_{\SS_X}\SS_X$, and so by the above identification we have
\[
\QF^{\vis}_{\xi}(\nu) = \QF^{\vis}_{\xi}(\Dual_{\SS_X}\SS_X) = \colim_{X}\QF^{\uni}(\SS)\otimes \xi ,
\]
so that the universal Poincaré form $\SS \to \QF^{\uni}(\SS)$ induces a map
\[
\colim_X \xi \to \QF^{\vis}_{\xi}(\nu) .
\]
Through this map for $\xi=\nu_X$, the Pontryagin-Thom collapse class $\SS\to \colim_X \nu_X$ provides us with the desired Poincaré form $q_{\mathrm{sig}}$ on $\nu_X$ for the $\nu_X$-twisted visible functor $\QF^\vis_{\nu_X}$. The Poincaré object $(\nu_X,q_{\mathrm{sig}}) \in \Poinc(\Spaf_X,\QF^{\vis}_{\nu_X})$ is the visible signature as constructed in loc.\ cit.


\begin{proposition}
\label{proposition:visible-signature-revisited}%
The visible signature $\vsign(X)$ of Definition~\reffour{definition:visible-signature} coincides with the visible signature $(\nu_X,q_{\mathrm{sig}})$ as defined in \S\refone{subsection:parametrised-spectra}.
In particular, under the embedding $e_{\Pro}$ the Poincaré object $(\nu_X,q_{\mathrm{sig}})$ maps to the constant diagram $\SS_X\colon X \to \Pro(\Spaf)$ with value $\SS$, with the Poincaré form $q \in (\QF^{\uni})^X(\SS_X) = \lim_X \QF^{\uni}(\SS)$ corresponding to the composite $\SS \to \QF^{\uni}(\SS) \to \lim_X \QF^{\uni}(\SS)$, where the first map is the universal Poincaré form.
\end{proposition}



\begin{proof}
By construction, $e_{\Pro}(\vsign(X))$ corresponds to the constant diagram $\SS_X\colon X \to \Pro(\Spaf)$ with value $\SS$ with the constant universal Poincaré form $\ovl{q}^{\uni} \in \lim_X \QF^{\uni}(\SS)$. Since $e_{\Pro}$ is a hermitian embedding by Lemma~\reffour{lemma:e_Pro_hermitian} it will suffice to show that this description also holds for the image of $(\nu_X,q_{\mathrm{sig}})$.
Now at the level of objects we have
\[
e_{\Pro}(\nu_X) = \SS_X
\]
by~\eqreffour{equation:ind-to-pro},
and the hermitian part of $e_{\Pro}$ at $\nu_X$ then implements the equivalence
\[
\QF^{\vis}_{\nu_X}(\nu_X) \simeq \colim_X\QF^{\uni}(\SS)\otimes \nu_X \xrightarrow{\simeq} \lim_X\QF^{\uni}(\SS) .
\]
It then suffices to recall that the constant class $\SS\to \lim_X \SS$ corresponds under this identification to the Pontryagin-Thom collapse class $\SS\to \colim_X \nu_X$ (by definition of the latter).
\end{proof}

